\documentclass[conference]{IEEEtran}
\IEEEoverridecommandlockouts

% ── Fonts & math (Times-like professional serif) ─────────────────────
\usepackage{amsmath,amssymb,amsfonts}
\usepackage{newtxtext,newtxmath}
\usepackage[T1]{fontenc}
\usepackage{microtype}

% ── Color system (project brand: navy primary, one orange accent) ────
\usepackage{xcolor}
\definecolor{ink}{HTML}{182230}      % heading ink
\definecolor{navy}{HTML}{1A4A73}     % primary
\definecolor{accent}{HTML}{E8590C}   % single accent
\definecolor{muted}{HTML}{5D6C7B}    % captions/footer
\definecolor{rulegray}{HTML}{B9C2CC} % hairlines

\usepackage{graphicx}
\usepackage{textcomp}
\usepackage{url}
\usepackage{booktabs}                % professional table rules
\usepackage{colortbl}
\usepackage{balance}                 % balanced last-page columns
\usepackage[font=small,labelfont={bf,color=navy},labelsep=period,
            justification=raggedright,singlelinecheck=false]{caption}
\usepackage{titlesec}                % heading hierarchy
\usepackage{enumitem}
\usepackage{fancyhdr}
\usepackage[hidelinks]{hyperref}

% ── Heading design: larger, colored, well-spaced ─────────────────────
\titleformat{\section}
  {\normalfont\large\bfseries\color{ink}}
  {\textcolor{accent}{\thesection.}}{0.55em}{}
\titlespacing*{\section}{0pt}{15pt plus 3pt minus 2pt}{8pt plus 1pt}
\titleformat{\subsection}
  {\normalfont\normalsize\bfseries\color{navy}}
  {\thesubsection.}{0.45em}{}
\titlespacing*{\subsection}{0pt}{11pt plus 2pt minus 1pt}{5pt}

% ── Enumerated/itemized lists: tighter, hanging indent ───────────────
\setlist[enumerate]{leftmargin=1.35em, itemsep=3pt plus1pt, topsep=4pt}
\setlist[itemize]{leftmargin=1.1em, itemsep=2.5pt, topsep=3pt}

% ── Airier text block: block paragraphs, no first-line indents ───────
\linespread{1.045}
\setlength{\parindent}{0pt}
\setlength{\parskip}{3.2pt plus 0.8pt}

% ── Footer: running title + page number on a hairline ────────────────
\pagestyle{fancy}
\fancyhf{}
\renewcommand{\headrulewidth}{0pt}
\renewcommand{\footrulewidth}{0.3pt}
\renewcommand{\footrule}{\hbox to\headwidth{\color{rulegray}\leaders\hrule height \footrulewidth\hfill}}
\fancyfoot[L]{\footnotesize\color{muted}Rules-Exact Conversational Tax Advisory}
\fancyfoot[R]{\footnotesize\color{ink}\thepage}
\fancypagestyle{plain}{\fancyhf{}\renewcommand{\headrulewidth}{0pt}%
  \renewcommand{\footrulewidth}{0.3pt}%
  \fancyfoot[L]{\footnotesize\color{muted}Rules-Exact Conversational Tax Advisory}%
  \fancyfoot[R]{\footnotesize\color{ink}\thepage}}

\newcommand{\thead}[1]{\textbf{#1}}

\newcommand{\Rs}{Rs.~}

\begin{document}

\title{Rules-Exact Conversational Tax Advisory:\\
A Deterministic Engine with a Closed-Loop\\Acceptance Intelligence Layer}

\author{\IEEEauthorblockN{Jayant (12406581)}
\IEEEauthorblockA{CSE472 : DEEP LEARNING FOR NATURAL LANGUAGE PROCESSING\\
Roll No.~RK4P24MZA06 \\
{\color{navy}\url{https://github.com/Jayant1711/Tax_Assist-system}}\\
November 2026}}

\maketitle
\thispagestyle{plain}

\begin{abstract}
\normalfont\normalsize
Tax numbers are easy to get wrong and hard for a reader to check, and
language models get them wrong fluently. This report describes Tax
Assist, a conversational advisor for Indian income tax that I built
around a single rule: nothing outside the tax calculator is allowed to
compute a rupee. The language layer only decides what the user
probably meant. A versioned calculator prices the old and new tax
regimes exactly, and every recommendation is ranked by re-running that
calculator on a hypothetical top-up, so an ``estimated savings''
figure is really the calculator agreeing with itself twice. Two
results from the work seem worth a reader's attention. I trained a
neuro-fuzzy (TSK) head to imitate the calculator and it got worse
every time I gave it more data; at 378k labeled rows its savings error
was off by a factor of roughly fifteen thousand, and the reason is now
clear, namely that the target contains a regime-flip $min()$ which
additive rule bases smooth over instead of pinning. And since a
confidence number nobody checks is just decoration, the system records
what users actually do with each recommendation and computes Expected
Calibration Error over those outcomes; a known-answer canary pins the
instrument to zero, and a standing gate refuses any learned confidence
until it measures ECE $\leq 0.05$ on at least thirty resolved events.
The shipped system passes 1{,}166 backend and 31 frontend tests, with
extraction certified at macro-F1 1.000. The falsified experiment also
reshaped the roadmap: the neuro-fuzzy layer is being repositioned not
as a regression model but as an \emph{explainable} rule base meant to
capture how a human chartered accountant reasons about savings, which
the exact calculator would then price.
\end{abstract}

\begin{IEEEkeywords}
\normalfont\normalsize
conversational AI, rule-based systems, neuro-fuzzy inference,
confidence calibration, negative results, information extraction,
tax computation
\end{IEEEkeywords}

\section{Introduction}
Ask a chatbot about Section 80C of the Indian Income-tax Act and you
will usually get a confident, well-formatted paragraph, sometimes even
with the right numbers. The difficulty is telling when it didn't.
Indian income tax is an unforgiving place to find out: liability comes
out of two parallel regimes interacting with over twenty deduction
sections, a senior-citizen matrix, surcharge with marginal relief, and
a rebate clause, and the person reading the answer generally has no
way to check any of it.

Early in this project a test profile earning Rs.~46 lakh produced a
report in which every deduction saved exactly nothing. I spent the
better part of a day convinced the calculator was broken. It wasn't.
The new regime really does disallow those deductions, and my interface
had simply failed to say so. That episode did more to shape the system
than any design document I wrote, because it made the goal concrete:
the fix was never ``trust the model less,'' it was ``show the
working, from the calculator, every time.''

Tax Assist is the system that came out of that. It is a web
application in which a user describes their income and investments the
way they would describe them to a human CA, in English or Hinglish,
and gets back a full computation plus a ranked set of opportunities.
The claim I defend is structural rather than statistical: no component
of the system can invent a rupee. I enforce that in one place (the
rules engine is the only calculator), verify it in another
(recommendations are priced by re-running the engine), and watch it
continuously (an append-only log and an event store that remembers
what users did with the advice).

Six things in this report deserve mention as contributions. The
two-pass pricing mechanism in Section~IV makes every savings figure a
recomputation rather than an estimate. The decision rules in
Section~V decide what to show and are exact by construction.
Section~V also contains what I consider the most useful result of the
whole project, a falsification: the ``obvious'' learned head failed at
three data scales, worse each time, and I can explain the mechanism
precisely. Section~VI describes the acceptance loop and the
calibration gate. Section~VII covers verification, where I want to
flag one habit that paid for itself repeatedly, which is treating a
five-minute manual session in the running product as a first-class
test. And Section~VIII states where I intend to take the neuro-fuzzy
layer next: away from regression, toward an explainable encoding of
CA-style savings reasoning.

\section{Background and System Overview}
\subsection{The tax domain, briefly}
A reader unfamiliar with Indian income tax needs three facts. First,
every taxpayer picks one of two regimes: the \emph{old regime}, which
allows deductions under numbered sections of the Act (80C for
investments up to Rs.~1.5 lakh, 80D for health insurance, 80E for
education-loan interest, 24(b) for home-loan interest, and so on) but
taxes income at higher slab rates; and the \emph{new regime}
(Section~115BAC), which offers lower slab rates and a larger standard
deduction but discards almost all of those deductions. Second, each
deduction section has a statutory cap, so every head has a remaining
\emph{headroom} that a recommendation can fill. Third, the interaction
is genuinely nonlinear: near certain income levels, topping up one
deduction \emph{flips} which regime is cheaper, so the saving from a
top-up can be discontinuous in the input. Section V returns to this
point, because it is exactly where a learned model of the savings
fails.

\subsection{The system}
Tax Assist is a three-tier web application. The client is Next.js and
React and presents a chat workspace plus the final advisory report.
The backend is FastAPI with seventeen endpoints covering
authentication, conversation, calculation, session management,
recommendation events, and funnel analytics. Requests pass through
wire-level validation, per-route rate limits, and an origin allow-list
before anything touches application logic. A database holds users,
server-side sessions, and the recommendation events; an append-only
audit log records every interaction.

I fixed four decisions early and never revisited them, stated here
plainly rather than by their internal names. The rules engine is the
only calculator in the system; anything that needs a number asks it.
The whole system runs offline and deterministically, so the same
session state produces the same numbers on any machine. Extraction
from free text stays rule-based rather than learned, a decision the
evaluation numbers support. And every recommendation shown to a user
is persisted as an event the user can later resolve with an action,
which is what makes the feedback half of this paper possible.

\section{Language Understanding}
\subsection{Processing chain}
Every utterance goes through the same chain. Prepass guards mask or
book percent signs, dates, and flat statutory amounts before any
parsing happens. An Indian-numeral parser then handles digit forms
(\texttt{1,25,000}, \texttt{18~LPA}, \texttt{45k~pm}), word forms
(``\emph{thirty-five lakhs}''), and Hinglish numerals
(``\emph{paanch hazaar}''). A longest-match guard keeps statutory
tokens from being misread; without it the section token
\texttt{80u\_80dd} decays into the amount ``80'' with a stray unit
``u,'' which turns a deduction section into eighty rupees of income.
Finally, semantic attention decides which category an amount belongs
to, and everything is clamped to a statutory range.

One correctness detail took me longer than it should have: clause
splitting. A comma inside a digit run is grouping, not grammar, and a
naive split on commas destroys \texttt{1,25,000} into three pieces.
The production rule splits on a comma or semicolon only when neither
neighbor is a digit, which as a regex is just two zero-width
assertions:
\begin{equation}
\texttt{(?<!\textbackslash d)[,;](?!\textbackslash d)}
\end{equation}

\subsection{Semantic attention}
Keywords vote for categories with a weight that decays by distance, so
an amount binds to the statutory context nearest to it:
\begin{equation}
s(k) = 100 \cdot e^{-d_k^{2} / 2\sigma^{2}},
\label{eq:attn}
\end{equation}
where $d_k$ is the character distance from keyword $k$ to the amount
and $\sigma=40$. In the sentence ``invested 1.5 lakhs in PPF \ldots
pay 2L rent,'' Section~80C wins the \Rs{}1.5L at $d{=}6$ and the rent
keyword wins the \Rs{}2L at $d{=}34$, which is what a human reader
would say too. When nothing scores above threshold the amount falls
back to a profile default, flagged as provisional; a default can
never overwrite a value the user actually booked.

\subsection{Evaluation}
I evaluate the language layer on template-disjoint splits of a
2{,}226-record labeled corpus, meaning test utterances never share a
template with anything in training. Production extraction reaches
macro precision, recall and F1 (the harmonic mean of the two,
averaged over categories) of 1.000 across 21 categories on the
434-record evaluation split, abstentions included (17/17 correct),
and clears the $\geq 95\%$ certification gate on every income head and
major section (Table~\ref{tab:nlp}). The research layer's intent
classifier, a bag-of-words linear model over 8 intents, sits at
96.41\%. A sequence-labeling network trained on only 51 templated
samples reaches 0.143 entity accuracy, nowhere near usable, so it
never entered the product. The training-set size is the limiting
factor here rather than the architecture, and widening the template
coverage is left as future work.

\begin{table}[t]
\caption{Language-layer evaluation (template-disjoint splits).}
\label{tab:nlp}
\centering
\footnotesize
\begin{tabular}{@{}lll@{}}
\toprule
\rowcolor{black!6}
\thead{Task} & \thead{Model} & \thead{Result} \\
\midrule
Extraction (21 cat.) & deterministic & \textbf{F1 1.000} (gate PASS) \\
Abstention & threshold & 17/17 correct \\
Intent (8 classes) & linear classifier & 96.41\% acc. \\
Entities (5 types) & BiLSTM, 51 samples & 0.143 acc.\ (off-path) \\
\bottomrule
\end{tabular}
\end{table}

\section{The Deterministic Pricing Core}
The calculator carries FY-2024--25 and FY-2025--26 as versioned rule
registries: slab rates, surcharge with marginal relief, the
health-and-education cess, the Section~87A rebate, the HRA exemption
under its statutory least-of-three test, and per-section caps.
Computation is a pure function of validated session state, so the same
session yields the same numbers anywhere. The ML corpus used later in
this paper is generated by this engine as well, from a fixed seed with
LF-normalized writes and content-addressed manifests, which means a
dataset can be rebuilt byte-for-byte if anyone doubts it.

\subsection{Two-pass recommendation pricing}
For a candidate deduction head $i$ with remaining headroom $\delta_i$,
the calculator prices the session as booked (pass one) and prices it
again with the head topped up (pass two):
\begin{equation}
s_i \;=\; T_{\min}(x) - T_{\min}(x + \delta_i \cdot \mathbf{e}_i),
\end{equation}
where $T_{\min}$ is whichever regime's total tax is cheaper. Since
both passes are the authoritative calculator, $s_i$ is not a guess
about what the calculator would say. It is the calculator saying it
twice.

\subsection{Exact per-deduction impact}
The final report also prices each booked deduction by leaving it out
and recomputing, in both regimes:
\begin{equation}
\Delta_i \;=\; T_{r}(x_{-i}) - T_{r}(x),
\label{eq:impact}
\end{equation}
with $r$ the recommended regime and $x_{-i}$ the session with head $i$
zeroed. The sign convention makes $\Delta_i$ positive exactly when the
item delivered relief. I should admit that the first version of this
equation subtracted the wrong way, and a \Rs{}1{,}30{,}000 benefit
showed up in the interface as negative. No amount of code review
caught it; clicking through a live profile did. That experience is why
manual end-to-end sessions count as tests in this project now. An item
is labeled ``applied'' when the regime's taxable income actually
changes without it, a recomputation test that correctly allows
let-out home-loan interest under the new regime while still refusing
self-occupied interest. In the live check, a profile with \Rs{}5L of
education-loan interest and \Rs{}2L of let-out home-loan interest
gives $\Delta_{80E} = \Rs{}1{,}30{,}000$ and $\Delta_{24b} =
\Rs{}41{,}600$, both of which match a hand-worked slab calculation; a
control profile under the new regime returns \Rs{}0 for Section~80C
and prints the statutory reason for it.

\section{Intelligence: Exact Rules and a Falsified Alternative}

\subsection{What the system decides, and how}
Three plain questions decide what the final report shows the user:
should this head be recommended at all, how much would filling it
save, and how confidently should the system say so. Each question is
answered from the recomputed quantities above, never from a model.
Because the answers come from the same closed-form expressions the
calculator uses, the rules measure zero error on a 10{,}995-row
validation set: opportunity precision and recall of 1.0, savings
error \Rs{}0.00. That is not a lucky split. The rules cannot be wrong
unless the calculator itself is wrong.

\subsection{The learned alternative, falsified at three scales}
The obvious research step is to train a model to imitate those
decisions, and I did, three times, each time with more data than the
last (Table~\ref{tab:tsk}). The model was a Takagi--Sugeno--Kang
neuro-fuzzy network, the standard adaptive fuzzy-inference design
\cite{takagi1985, jang1993}, trained on engine-generated labels. The
result was failure in every configuration, getting reliably worse as
data grew. At the final scale the savings error reached 1{,}523{,}679\%
mean-percentage error, and the opportunity head never produced a
single useful prediction at any scale.

\begin{table}[t]
\caption{A learned neuro-fuzzy head across three data scales. More
data made every gate worse; the exact rules pass all gates by
construction.}
\label{tab:tsk}
\centering
\footnotesize
\begin{tabular}{@{}lrrr@{}}
\toprule
\rowcolor{black!6}
\thead{Scale} & \thead{Positives} & \thead{Savings error} & \thead{Opportunity} \\
\midrule
3.5k rows & 146 & fail & P 0 / R 0 \\
77k rows & 1{,}623 & MAPE 795\% & P 0 / R 0 \\
378k rows & 51{,}873 & \textbf{MAPE 1{,}523{,}679\%} & P 0 / R 0 \\
\bottomrule
\end{tabular}
\end{table}

It took a while to accept that this was structural and not a tuning
problem on my part. The savings label is exactly linear in two ratios
the calculator can produce itself; when I hand the model those ratios
as features, validation error drops to about $1.6 \times 10^{-14}$, in
other words the label is a formula. The catch sits one level down: the
regime choice inside those ratios is a discontinuous $min()$, and an
additive rule base interpolates across the flip instead of pinning
it. The boundary signal gets diluted, which shows up as a ranking
accuracy of 65.26\% at the final scale against 100\% for the exact
rules. The opportunity label is even starker: a top-up helps exactly
when the current regime ignores that deduction, a condition that holds
without exception across all 377{,}753 rows, so there is a
distribution to learn only if one believes a rule is a distribution.
My conclusion, now written down as policy so I do not waste a future
month rediscovering it, is that learned regression on closed-form
targets is inadmissible no matter how much data arrives. The dataset
itself, 21{,}964 scenarios and 51{,}873 positive rows at 13.73\%
density, stays around as a deterministic benchmark rather than a
growth axis.

\section{Closing the Loop: Acceptance and Calibration}

\subsection{Acceptance events}
Every ranked recommendation shown at the final report is written to a
store as a shown-event carrying the deduction head, the top-up, the
exact savings, and the confidence. The user resolves it later with an
action: accepted, rejected, ignored, viewed, or dismissed. A funnel
per session aggregates totals and per-head adoption, which is what
lets the system eventually learn which kinds of advice people actually
take.

Two implementation choices turned out to matter more than they look.
Event identifiers are derived deterministically from the session and
the head, so the identifier the interface renders is the identifier in
the store. My first version generated random identifiers per call,
which meant every card click would have hit an event that did not
exist; I found that by clicking, not by reading code. Second, the
closing-turn intelligence lives in one shared code path used by both
the natural conversational close and the ``\texttt{/end}'' override.
The override originally returned early and skipped the intelligence
entirely, so sessions closed that way recorded no events at all. Both
mistakes are now pinned by regression tests. In the live end-to-end
check, a session shows three recommendations (Sections 24B, 80C and
80G, each at confidence 0.85), the user accepts one, and the funnel
reports 3 shown, 1 accepted, rate 0.333, with per-head exact savings.

\subsection{Calibration and the standing gate}
The confidence the system serves today is a constant: 0.85 when the
calculator verified the number, 0.30 for priced-no-benefit. A constant
is defensible precisely because its meaning can be stated in one line.
Whether a learned confidence could replace it is a question I chose to
instrument rather than guess at. Expected Calibration Error over
equal-width bins,
\begin{equation}
\mathrm{ECE} = \sum_{b=1}^{B} \frac{|B_b|}{N}
\left| \mathrm{acc}(B_b) - \mathrm{conf}(B_b) \right|,
\label{eq:ece}
\end{equation}
is computed over resolved acceptance events. The instrument itself is
validated by a known-answer canary: I construct 200 pairs whose
observed acceptance equals each bin's average confidence exactly, and
the implementation must report ECE 0.0 on them and must flag the
flipped outcome set as miscalibrated. It does both, but the first
canary did not; with ten items per bin, Python's banker's rounding
turned a perfect construction into a reported ECE of 0.05, which is
the sort of bug that would have quietly discredited the whole
instrument. Twenty items per bin makes every bin's positive count an
exact integer and the artifact disappears. The standing gate is then
simple: no learned confidence ships until it measures ECE
$\leq 0.05$ on at least 30 resolved events without hurting the
acceptance rate. The store currently holds one resolved event, so the
honest measurement today is ``unmeasured,'' and the anchor stays.

\section{Verification}
Table~\ref{tab:ver} is the verification picture. Backend tests cover
the calculator against hand-computed statutory cases, the language
chain including adversarial numbering and Hinglish, session ownership
and rate limits, the event store and its action endpoints, and the
impact calculation's sign and new-regime behavior, that last pair
pinned after the live session caught the defects. Continuous
integration runs the backend suite, the frontend suite, and scripted
browser end-to-end flows on every push. Two tagged releases went out
through CI-gated pull requests, and I re-ran the full suite on the
merged result both times.

\begin{table}[t]
\caption{Verification summary of the shipped system.}
\label{tab:ver}
\centering
\footnotesize
\begin{tabular}{@{}ll@{}}
\toprule
\rowcolor{black!6}
\thead{Suite} & \thead{Result} \\
\midrule
Backend (pytest) & 1{,}166 passed / 0 failed \\
Frontend (vitest + tsc) & 31 passed / clean \\
End-to-end (Playwright) & 7 flows green \\
Extraction gate & macro F1 1.000 (PASS) \\
Decision rules & P/R 1.0, savings MAE \Rs{}0.00 \\
Calibration canary & ECE 0.0 (framework exact) \\
CI & 3 jobs green per push \\
\bottomrule
\end{tabular}
\end{table}

\section{Discussion: What the Neuro-Fuzzy Layer Is For}
I want to close the technical part by being precise about where a
fuzzy system genuinely belongs in this design, because the failed
experiment above is easy to misread as ``neuro-fuzzy does not work
here,'' and that is not the lesson I take.

A neuro-fuzzy inference system is, structurally, a set of readable
rules. Each rule reads roughly like: \emph{IF income is high AND 80C
headroom is large AND the taxpayer is young, THEN an NPS top-up is a
strong opportunity}. The fuzzy membership functions give the
conditions their graded, human flavor --- ``high,'' ``large,'' ``near
the rebate cliff'' --- and the rule base as a whole is inspectable:
unlike a neural network's weights, a trained fuzzy system can be
translated back into sentences a person can argue with. That is exactly
the shape of how a chartered accountant reasons about savings. A good
CA does not regress a number; they carry a set of graded, conditional
judgments --- who is this client, which levers are still open, what is
the interaction with the regime choice --- and combine them into a
plan.

So the direction I am taking the neuro-fuzzy layer is not regression
but \emph{representation}. Concretely, three steps. First, make the
rule base explainable: every rule that fires for a given client should
be translatable into one plain sentence of CA-style reasoning, shown
alongside the exact numbers. Second, validate the decoded rules
against a practicing CA's stated reasoning on the same profiles, so
the fuzzy layer is judged as a model of professional judgment, not
just as a predictor. Third, once the rules are explainable and
validated, let the layer propose \emph{plans} rather than single
top-ups --- combinations like ``add a Rs.~50k NPS contribution, move
to the old regime, and claim the pending 80D headroom'' --- with the
deterministic calculator pricing every proposed plan exactly and the
calibration gate governing how confidently each step is phrased. The
engine prices; the fuzzy layer learns, in readable form, to think like
a CA; and nothing either of them says is trusted until the engine
confirms the rupees. That division of labor is the honest lesson of
the falsification, turned into a research program instead of a dead
end.

\section{Limitations}
Salaried profiles are where the advisor is deep; business and
capital-gains flows exist but are thinner, and I would not trust the
business paths for anything real yet. FY~2025--26 is the only certified
year, and a new fiscal year means adding a rules registry and
re-certifying, which the design anticipates but which has not been
exercised. The calibration gate is built and proven as an instrument
but currently uninformative, since one resolved event is not a sample.
The explainable-CA direction in Section~VIII is a design intent, not a
result; nothing in this paper should be read as claiming the fuzzy
layer currently reasons like a CA. Personalization from acceptance
data is deliberately blocked until the store holds roughly fifty
resolved events. And the system advises; it does not file anything
with anyone.

\section{Conclusion and Future Work}
The lesson I would keep from this project is that a tax advisor can be
exact, and that exactness is an architectural property rather than a
model quality. One deterministic calculator, everything else
subordinate, costs almost nothing in capability, and it converts a
trust problem into an engineering one. Every figure is recomputable,
every claim ships with its evidence, and every failed idea is recorded
alongside the measurement that killed it. If the project has a single
most valuable machine-learning artifact, it is a negative result with
the mechanism written down; its most valuable gate is one that keeps a
model out until the data exists to judge it.

What comes next follows the evidence. An effectiveness dashboard over
the acceptance funnel, with the Accepted-Rupee metric as the headline.
One-click application of a recommendation through the calculator's own
write path. And the neuro-fuzzy layer rebuilt toward the goal it was
always meant for: an explainable rule base that captures how a
chartered accountant thinks about legal savings, validated against a
practitioner's judgment, proposing full savings plans that the exact
calculator prices line by line. A FY~2026--27 rules drop will be the
first real test of the registry's forward compatibility.

\balance
\begin{thebibliography}{12}

\bibitem{takagi1985}
T.~Takagi and M.~Sugeno, ``Fuzzy identification of systems and its
applications to modeling and control,'' \emph{IEEE Trans. Syst., Man,
Cybern.}, vol.~SMC-15, no.~1, pp. 116--132, 1985.

\bibitem{jang1993}
J.-S.~R. Jang, ``ANFIS: Adaptive-network-based fuzzy inference
system,'' \emph{IEEE Trans. Syst., Man, Cybern.}, vol.~23, no.~3,
pp. 665--685, 1993.

\bibitem{guo2017}
C.~Guo, G.~Pleiss, Y.~Sun, and K.~Q. Weinberger, ``On calibration of
modern neural networks,'' in \emph{Proc. 34th Int. Conf. Machine
Learning (ICML)}, 2017, pp. 1321--1330.

\bibitem{chiticariu2013}
L.~Chiticariu, Y.~Li, and F.~R. Reiss, ``Rule-based information
extraction revisited,'' in \emph{Proc. 51st Annu. Meeting Assoc.
Computational Linguistics (ACL)}, 2013, pp. 406--416.

\bibitem{bahdanau2015}
D.~Bahdanau, K.~Cho, and Y.~Bengio, ``Neural machine translation by
jointly learning to align and translate,'' in \emph{Proc. 3rd Int.
Conf. Learning Representations (ICLR)}, 2015.

\bibitem{luong2015}
M.-T. Luong, H.~Pham, and C.~D. Manning, ``Effective approaches to
attention-based neural machine translation,'' in \emph{Proc. EMNLP},
2015, pp. 1412--1421.

\bibitem{sutskever2014}
I.~Sutskever, O.~Vinyals, and Q.~V. Le, ``Sequence to sequence
learning with neural networks,'' in \emph{Proc. NeurIPS}, 2014,
pp. 3104--3112.

\bibitem{mikolov2013}
T.~Mikolov, K.~Chen, G.~Corrado, and J.~Dean, ``Efficient estimation
of word representations in vector space,'' in \emph{Proc. ICLR
Workshop}, 2013.

\bibitem{pennington2014}
J.~Pennington, R.~Socher, and C.~D. Manning, ``GloVe: Global vectors
for word representation,'' in \emph{Proc. EMNLP}, 2014,
pp. 1532--1543.

\bibitem{devlin2019}
J.~Devlin, M.-W. Chang, K.~Lee, and K.~Toutanova, ``BERT:
Pre-training of deep bidirectional transformers for language
understanding,'' in \emph{Proc. NAACL-HLT}, 2019, pp. 4171--4186.

\bibitem{vaswani2017}
A.~Vaswani \emph{et~al.}, ``Attention is all you need,'' in
\emph{Proc. NeurIPS}, 2017, pp. 5998--6008.

\bibitem{ita1961}
\emph{Income-tax Act, 1961} (India), \S115BAC, inserted by the
Finance Act, 2020; scheme of taxation under the new regime.

\end{thebibliography}

\end{document}
